% GENERATED FILE. Edit canonical lessons, then run npm run generate:books.
% canonical-source-hash: fnv1a64:9c12c869c9e5b661
% canonical-lessons: MR-C406-padpath, MR-C406-padcari, MR-C406-phalak, MR-C406-uddanpul, MR-C406-nyahalne

\chapter{Walking in Town}
\label{ch:mr-walking-in-town}
\hlchaptermodality{406}

\begin{chapteropening}
\textbf{By the end of this chapter:} \emph{I can talk about a pavement, pedestrians, signboards and a flyover.}

Complete the last lesson of chapter 406: I can talk about a pavement, pedestrians, signboards and a flyover.
\end{chapteropening}

\section[\texorpdfstring{\mr{पदपथ} (padpath)}{पदपथ (padpath)}]{\mr{पदपथ} (padpath) — a pavement; a footpath}

\label{lesson:MR-C406-padpath}



\begin{quote}
Before the new one: say the Marathi for border customs duty, then the Marathi for a visa.
\end{quote}

\subsection*{You'll want to know: \mr{पदपथ}}
\textbf{\mr{पदपथ}} — \emph{padpath} — \textquotedblleft{}a pavement; a footpath\textquotedblright{}.

\begin{quote}
\mr{या} \mr{पदपथावर} \mr{भाजीवाले} \mr{बसतात}.
\end{quote}

\begin{quote}
\emph{yā padpathāvar bhājīvāle bastāt.}
\end{quote}

\begin{quote}
Vegetable sellers sit on this pavement.
\end{quote}

Masculine.

A formal word for the raised walk beside a town road.

\begin{cousinweb}
\textbf{\mr{पदपथ}} (\emph{padpath}) is built from Sanskrit \emph{pada}, \textquotedblleft{}foot, step\textquotedblright{}, and \emph{patha}, \textquotedblleft{}path\textquotedblright{} (Monier-Williams): a path for feet.
\end{cousinweb}

\begin{grammarlens}[title={What you've built}]
One more word for getting around.
\end{grammarlens}

\subsection*{Your turn}
\begin{itemize}
\raggedright
  \item \emph{Say it:} \emph{padpath}
  \item \emph{Say it:} \emph{padpath}, once more
  \item \emph{Say it:} \emph{yā padpathāvar bhājīvāle bastāt.} — Vegetable sellers sit on this pavement.
  \item \emph{Recall:} say the Marathi for border customs duty, then the Marathi for a visa, then say \emph{padpath} again
\end{itemize}

\begin{tcolorbox}[breakable,colback=teal!4,colframe=teal!35!black,title={Before you move on}]
What does \mr{पदपथ} mean? (\textquotedblleft{}A pavement; a footpath\textquotedblright{}.) Say it once more.
\end{tcolorbox}
\section[\texorpdfstring{\mr{पादचारी} (pādcārī)}{पादचारी (pādcārī)}]{\mr{पादचारी} (pādcārī) — a pedestrian}

\label{lesson:MR-C406-padcari}



\begin{quote}
Before the new one: say the Marathi for a visa, then the Marathi for a pavement.
\end{quote}

\subsection*{You'll want to know: \mr{पादचारी}}
\textbf{\mr{पादचारी}} — \emph{pādcārī} — \textquotedblleft{}a pedestrian\textquotedblright{}.

\begin{quote}
\mr{पादचारी} \mr{पदपथावरून} \mr{चालतात}.
\end{quote}

\begin{quote}
\emph{pādcārī padpathāvarūn cāltāt.}
\end{quote}

\begin{quote}
Pedestrians walk on the pavement.
\end{quote}

Masculine.

The plural is the same: \mr{एक} \mr{पादचारी}, \mr{खूप} \mr{पादचारी}.

\begin{cousinweb}
\textbf{\mr{पादचारी}} (\emph{pādcārī}) is a tatsama, taken from Sanskrit \emph{pādacārin}, \textquotedblleft{}going on foot, walking\textquotedblright{} (Monier-Williams).
\end{cousinweb}

\begin{grammarlens}[title={What you've built}]
One more word for getting around.
\end{grammarlens}

\subsection*{Your turn}
\begin{itemize}
\raggedright
  \item \emph{Say it:} \emph{pādcārī}
  \item \emph{Say it:} \emph{pādcārī}, once more
  \item \emph{Say it:} \emph{pādcārī padpathāvarūn cāltāt.} — Pedestrians walk on the pavement.
  \item \emph{Recall:} say the Marathi for a visa, then the Marathi for a pavement, then say \emph{pādcārī} again
\end{itemize}

\begin{tcolorbox}[breakable,colback=teal!4,colframe=teal!35!black,title={Before you move on}]
What does \mr{पादचारी} mean? (\textquotedblleft{}A pedestrian\textquotedblright{}.) Say it once more.
\end{tcolorbox}
\section[\texorpdfstring{\mr{फलक} (phalak)}{फलक (phalak)}]{\mr{फलक} (phalak) — a signboard; a notice board}

\label{lesson:MR-C406-phalak}



\begin{quote}
Before the new one: say the Marathi for a pavement, then the Marathi for a pedestrian.
\end{quote}

\subsection*{You'll want to know: \mr{फलक}}
\textbf{\mr{फलक}} — \emph{phalak} — \textquotedblleft{}a signboard; a notice board\textquotedblright{}.

\begin{quote}
\mr{दुकानावर} \mr{नवीन} \mr{फलक} \mr{लावला} \mr{आहे}.
\end{quote}

\begin{quote}
\emph{dukānāvar navīn phalak lāvlā āhe.}
\end{quote}

\begin{quote}
A new signboard has been put up on the shop.
\end{quote}

Masculine.

\mr{फलक} \mr{लावणे}, to put up a sign.

\begin{cousinweb}
\textbf{\mr{फलक}} (\emph{phalak}) is a tatsama, taken from Sanskrit \emph{phalaka}, \textquotedblleft{}a board, a plank\textquotedblright{} (Monier-Williams).
\end{cousinweb}

\begin{grammarlens}[title={What you've built}]
One more word for getting around.
\end{grammarlens}

\subsection*{Your turn}
\begin{itemize}
\raggedright
  \item \emph{Say it:} \emph{phalak}
  \item \emph{Say it:} \emph{phalak}, once more
  \item \emph{Say it:} \emph{dukānāvar navīn phalak lāvlā āhe.} — A new signboard has been put up on the shop.
  \item \emph{Recall:} say the Marathi for a pavement, then the Marathi for a pedestrian, then say \emph{phalak} again
\end{itemize}

\begin{tcolorbox}[breakable,colback=teal!4,colframe=teal!35!black,title={Before you move on}]
What does \mr{फलक} mean? (\textquotedblleft{}A signboard; a notice board\textquotedblright{}.) Say it once more.
\end{tcolorbox}
\section[\texorpdfstring{\mr{उड्डाणपूल} (uḍḍāṇpūl)}{उड्डाणपूल (uḍḍāṇpūl)}]{\mr{उड्डाणपूल} (uḍḍāṇpūl) — a flyover}

\label{lesson:MR-C406-uddanpul}



\begin{quote}
Before the new one: say the Marathi for a pedestrian, then the Marathi for a signboard.
\end{quote}

\subsection*{You'll want to know: \mr{उड्डाणपूल}}
\textbf{\mr{उड्डाणपूल}} — \emph{uḍḍāṇpūl} — \textquotedblleft{}a flyover\textquotedblright{}.

\begin{quote}
\mr{उड्डाणपुलामुळे} \mr{चौकातली} \mr{गर्दी} \mr{कमी} \mr{झाली}.
\end{quote}

\begin{quote}
\emph{uḍḍāṇpulāmuḷe caukātlī gardī kamī jhālī.}
\end{quote}

\begin{quote}
Thanks to the flyover the crowding at the square went down.
\end{quote}

Masculine.

Before an ending the long \mr{ऊ} of \mr{पूल} shortens: \mr{उड्डाणपुलावर}.

\begin{cousinweb}
\textbf{\mr{उड्डाणपूल}} (\emph{uḍḍāṇpūl}) is a Marathi compound of \textbf{\mr{उड्डाण}}, \textquotedblleft{}a flight, a leap up\textquotedblright{}, built on Sanskrit \emph{uḍḍayana}, \textquotedblleft{}flying up\textquotedblright{} (Monier-Williams), and \textbf{\mr{पूल}}, \textquotedblleft{}a bridge\textquotedblright{}: a bridge that leaps over a road.
\end{cousinweb}

\begin{grammarlens}[title={What you've built}]
One more word for getting around.
\end{grammarlens}

\subsection*{Your turn}
\begin{itemize}
\raggedright
  \item \emph{Say it:} \emph{uḍḍāṇpūl}
  \item \emph{Say it:} \emph{uḍḍāṇpūl}, once more
  \item \emph{Say it:} \emph{uḍḍāṇpulāmuḷe caukātlī gardī kamī jhālī.} — Thanks to the flyover the crowding at the square went down.
  \item \emph{Recall:} say the Marathi for a pedestrian, then the Marathi for a signboard, then say \emph{uḍḍāṇpūl} again
\end{itemize}

\begin{tcolorbox}[breakable,colback=teal!4,colframe=teal!35!black,title={Before you move on}]
What does \mr{उड्डाणपूल} mean? (\textquotedblleft{}A flyover\textquotedblright{}.) Say it once more.
\end{tcolorbox}
\section[\texorpdfstring{\mr{न्याहाळणे} (nyāhāḷṇe)}{न्याहाळणे (nyāhāḷṇe)}]{\mr{न्याहाळणे} (nyāhāḷṇe) — to gaze at}

\label{lesson:MR-C406-nyahalne}



\begin{quote}
Before the new one: say the Marathi for a signboard, then the Marathi for a flyover.
\end{quote}

\subsection*{You'll want to know: \mr{न्याहाळणे}}
\textbf{\mr{न्याहाळणे}} — \emph{nyāhāḷṇe} — \textquotedblleft{}to gaze at\textquotedblright{}.

\begin{quote}
\mr{मुलं} \mr{दुकानातली} \mr{खेळणी} \mr{न्याहाळत} \mr{होती}.
\end{quote}

\begin{quote}
\emph{mulaṁ dukānātlī kheḷṇī nyāhāḷat hotī.}
\end{quote}

\begin{quote}
The children were gazing at the toys in the shop.
\end{quote}

Past, with -\mr{ने} (he …): \mr{त्याने} \mr{न्याहाळलं}, \emph{tyāne nyāhāḷlaṁ}.

A slow, careful look; the thing looked at is the object.

\begin{cousinweb}
\textbf{\mr{न्याहाळणे}} (\emph{nyāhāḷṇe}) is a Marathi verb for looking at something slowly and closely; its origin is not settled.
\end{cousinweb}

\begin{grammarlens}[title={What you've built}]
One more word for getting around.
\end{grammarlens}

\subsection*{Your turn}
\begin{itemize}
\raggedright
  \item \emph{Say it:} \emph{nyāhāḷṇe}
  \item \emph{Say it:} \emph{nyāhāḷṇe}, once more
  \item \emph{Say it:} \emph{mulaṁ dukānātlī kheḷṇī nyāhāḷat hotī.} — The children were gazing at the toys in the shop.
  \item \emph{Recall:} say the Marathi for a signboard, then the Marathi for a flyover, then say \emph{nyāhāḷṇe} again
\end{itemize}

\begin{tcolorbox}[breakable,colback=teal!4,colframe=teal!35!black,title={Before you move on}]
What does \mr{न्याहाळणे} mean? (\textquotedblleft{}To gaze at\textquotedblright{}.) Say it once more.
\end{tcolorbox}
